\section{Grup $\mathbf{\emph{K}_0}$ untuk Aljabar-$C^*$ Beridentitas}

\begin{defn} Misalkan $A$ aljabar-$C^*$ beridentitas. Misalkan
 \index{K0@$K_0(A)$!untuk $A$ beridentitas}%
$K_0(A) := G(\fml D(A))$, grup Grothendieck dari semigrup $\fml D(A)$ yang didefinisikan dalam~\ref{0060411} dan~\ref{0060414},
dan definisikan
 \index{<bracsqr@$[p\,]$ (citra $[p\,]_{\fml D}$ di bawah pemetaan Grothendieck)}%
   \[ [\hphantom{P}]\colon \fml P_\infty(A) \sto K_0(A)
           \colon p \mapsto \sbsb{\gamma}{\fml D(A)}\bigl([p\,]_{\fml D}\bigr)\,. \]
\end{defn}

Proposisi berikut merupakan konsekuensi langsung dari definisi ini.

\begin{prop}\label{0062188fa} Misalkan $A$ aljabar-$C^*$ beridentitas dan $p$, $q \in \fml P_\infty(A)$. Jika $p$ dan $q$
ekuivalen Murray--von Neumann, maka $[p\,] = [q\,]$ dalam~$K_0(A)$.
\end{prop}

\begin{defn} Misalkan $A$ aljabar-$C^*$ beridentitas dan $p$, $q \in \fml P_\infty(A)$. Kita mengatakan bahwa $p$
 \index{ekuivalensi!stabil}%
 \index{ekuivalensi stabil}%
\df{ekuivalen secara stabil} dengan $q$ dan menulis
 \index{<binrelequivst@$p \sim_{st} q$ (ekuivalensi stabil)}%
$p \sim_{st} q$ jika terdapat proyeksi $r \in \fml P_\infty(A)$ sedemikian sehingga $p \oplus r \sim q \oplus r$.
\end{defn}

\begin{prop} Ekuivalensi stabil $\sim_{st}$ merupakan relasi ekuivalensi pada $\fml P_\infty(A)$.
\end{prop}

\begin{prop} Misalkan $A$ aljabar-$C^*$ beridentitas dan $p$, $q \in \fml P_\infty(A)$. Maka $p
\sim_{st} q$ jika dan hanya jika $p \oplus \vc 1_n \sim q \oplus \vc 1_n$ untuk suatu $n \in \N$.
\end{prop}

Di sini, tentu saja, $\vc 1_n$ adalah identitas perkalian dalam $\M nA$.

\begin{prop}[Gambaran Standar $K_0(A)$ ketika $A$ beridentitas]\label{0062224} Jika $A$ suatu
aljabar-$C^*$ beridentitas, maka
 \index{standar!gambaran $K_0(A)$!ketika $A$ beridentitas}%
 \index{gambaran (\seeonly{gambaran standar}))}%
   \begin{align*}
         K_0(A) &= \{ [p\,] - [q\,]\colon p,q \in \fml P_\infty(A)\}  \\
                &= \{ [p\,] - [q\,]\colon  p,q \in \fml P_n(A) \text{ untuk suatu $n \in \N$}\}\,.
   \end{align*}
\end{prop}

\begin{prop} Misalkan $A$ aljabar-$C^*$ beridentitas dan $p$, $q \in \fml P_\infty(A)$. Maka $[p
\oplus q\,] = [p\,] + [q\,]$.
\end{prop}

\begin{prop} Misalkan $A$ aljabar-$C^*$ beridentitas dan $p$, $q \in \fml P_n(A)$. Jika $p \sim_h q$
dalam $\fml P_n(A)$, maka $[p\,] = [q\,]$.
\end{prop}

\begin{prop} Misalkan $A$ aljabar-$C^*$ beridentitas dan $p$, $q \in \fml P_n(A)$. Jika $p \perp q$
dalam $\fml P_n(A)$, maka $p + q \in \fml P_n(A)$ dan $[p + q\,] = [p\,] + [q\,]$.
\end{prop}

\begin{prop} Misalkan $A$ aljabar-$C^*$ beridentitas dan $p$, $q \in \fml P_\infty(A)$. Maka
$[p\,] = [q\,]$ jika dan hanya jika $p \sim_{st} q$.
\end{prop}

 Proposisi berikut menetapkan suatu sifat universal dari $K_0(A)$ ketika $A$ beridentitas. Di dalamnya, funktor pelupa $\abs{\hphantom{G}}$
 adalah funktor yang dijelaskan dalam proposisi~\ref{0062151}.

\begin{prop}\label{0062244} Misalkan $A$ aljabar-$C^*$ beridentitas, $G$ grup Abelian, dan $\nu\colon \fml P_\infty(A) \sto \abs G$
homomorfisme semigrup yang memenuhi
 \index{universal!sifat!dari $K_0$ (kasus beridentitas)}%
 \index{K0@$K_0$!sifat universal dari}%
 \begin{enumerate}[label=\rm{(\alph*)}]
    \item $\nu(\vc 0_A) = \vc 0_G$ dan
    \item jika $p \sim_h q$ dalam $\fml P_n(A)$ untuk suatu $n$, maka $\nu(p) = \nu(q)$.
 \end{enumerate}
Maka terdapat homomorfisme grup tunggal $\wt\nu\colon K_0(A) \sto G$ sedemikian sehingga diagram berikut komutatif.
   \[\xy
      \qtriangle/{>}`{>}`{-->}/<1000,700>[\fml P_\infty(A)`\abs{K_0(A)}`\abs G;{[\hphantom{P}]}`\nu`\abs{\wt\nu}]
      \morphism(1600,700)|r|/{-->}/<0,-700>[K_0(A)`G;\wt\nu]
     \endxy\]
\end{prop}

\begin{proof}[\emph{Petunjuk untuk bukti}] Misalkan $\tau\colon \fml D(A) \sto \abs G\colon \sbsb{[p\,]}{\fml D} \mapsto \nu(p)$. % SOURCE-CORRECTION: CH17-C026
Periksa bahwa $\tau$ terdefinisi dengan baik dan merupakan homomorfisme semigrup. Kemudian gunakan proposisi~\ref{0062151}.  \ns
\end{proof}

\begin{defn} Homomorfisme-$*\,$ $\phi\colon A \sto B$ di antara aljabar-$C^*$ diperluas, untuk
setiap $n \in \N$, menjadi homomorfisme-$*\,$ $\phi\colon \M nA \sto \M nB$ dan juga (karena homomorfisme-$*\,$ membawa proyeksi
ke proyeksi) memiliki pembatasan berupa pemetaan $\phi$ dari $\fml P_\infty(A)$ ke $\fml P_\infty(B)$. % SOURCE-CORRECTION: CH17-C012
Untuk homomorfisme-$*\,$ $\phi$ semacam itu, definisikan
   \[ \nu\colon \fml P_\infty(A) \sto K_0(B)\colon p \mapsto [\phi(p)]\,. \]
Maka $\nu$ merupakan homomorfisme semigrup yang memenuhi syarat (a) dan (b) dalam proposisi~\ref{0062244}, sehingga
terdapat homomorfisme grup tunggal
 \index{K0@$K_0(\phi)$!kasus beridentitas}%
$K_0(\phi)\colon K_0(A) \sto K_0(B)$ sedemikian sehingga $K_0(\phi)\bigl([p\,]\bigr) = \nu(p)$ untuk setiap $p \in \fml P_\infty(A)$.
\end{defn}

\begin{prop} Pasangan pemetaan $A \mapsto K_0(A)$, $\phi \mapsto K_0(\phi)$ merupakan funktor kovarian
dari kategori aljabar-$C^*$ beridentitas dan homomorfisme-$*\,$ ke kategori grup Abelian dan
homomorfisme grup. Lebih lanjut, untuk semua homomorfisme-$*\,$ $\phi\colon A \sto B$ di antara aljabar-$C^*$ beridentitas, diagram
berikut komutatif.
  \[\xy
         \Square[\fml P_\infty(A)`\fml
         P_\infty(B)`K_0(A)`K_0(B);\phi`{[\hphantom{P}]}`{[\hphantom{P}]}`K_0(\phi)]
  \endxy\]
\end{prop}

\begin{notn} Untuk aljabar-$C^*$ $A$ dan $B$, misalkan
 \index{hom@$\Hom(A,B)$ (himpunan homomorfisme-$*\,$ di antara aljabar-$C^*$)}%
$\Hom(A,B)$ adalah keluarga semua homomorfisme-$*\,$ dari $A$ ke~$B$.
\end{notn}

\begin{defn} Misalkan $A$ dan $B$ aljabar-$C^*$ serta $a \in A$. Untuk $\phi$, $\psi \in \Hom(A,B)$,
misalkan
   \[ d_a(\phi,\psi) = \norm{\phi(a) - \psi(a)}\,. \]
Maka $d_a$ merupakan pseudometrik pada $\Hom (A,B)$. Topologi yang dibangkitkan oleh keluarga $\{d_a\colon a \in A\}$ adalah
 \index{topologi norma-titik}%
 \index{topologi!norma-titik}%
\df{topologi norma-titik} pada~$\Hom(A,B)$. (Untuk setiap $a \in A$, misalkan $\sfml B_a$ adalah keluarga bola terbuka dalam $\Hom(A,B)$ yang dibangkitkan oleh
pseudometrik~$d_a$. Keluarga $\bigcup\{\sfml B_a\colon a \in A\}$ merupakan subbasis bagi topologi norma-titik.)
\end{defn}

\begin{defn} Misalkan $A$ dan $B$ aljabar-$C^*$. Kita mengatakan bahwa homomorfisme-$*\,$ $\phi$, $\psi\colon A \sto B$
 \index{homotop!homomorfisme-$*\,$}%
 \index{<star@homomorfisme-$*\,$!homotopi dari}%
 \index{<binrelequivh@$\phi\sim_h\psi$ (homotopi homomorfisme-$*\,$)}%
\df{homotop}, dan menulis $\phi \sim_h \psi$, jika terdapat suatu fungsi
   \[ c\colon [0,1] \times A \sto B\colon (t,a) \mapsto c_t(a) \]
sedemikian sehingga
  \begin{enumerate}[label=\rm{(\alph*)}]
     \item untuk setiap $t \in [0,1]$, pemetaan $c_t\colon A \sto B \colon a \mapsto c_t(a)$ merupakan homomorfisme-$*\,$,
     \item untuk setiap $a \in A$, pemetaan $c(a)\colon [0,1] \sto B\colon t \mapsto c_t(a)$ merupakan lintasan (kontinu) dalam~$B$,
     \item $c_0 = \phi$, dan
     \item $c_1 = \psi$.
  \end{enumerate}
\end{defn}

\begin{prop} Dua homomorfisme-$*\,$ $\phi_0$, $\phi_1\colon A \sto B$ di antara aljabar-$C^*$
homotop jika dan hanya jika terdapat lintasan kontinu norma-titik dari $\phi_0$ ke $\phi_1$ dalam~$\Hom(A,B)$.
\end{prop}

\begin{defn} Kita mengatakan bahwa aljabar-$C^*$ $A$ dan $B$
 \index{homotop!ekuivalensi!homomorfisme-$*\,$}%
 \index{ekuivalensi!homotopi!homomorfisme-$*\,$}%
\df{ekuivalen secara homotopi} jika terdapat homomorfisme-$*\,$ $\phi\colon A \sto B$ dan $\psi\colon B \sto A$ sedemikian sehingga $\psi
\circ \phi \sim_h \id A$ dan $\phi \circ \psi \sim_h \id B$. Suatu aljabar-$C^*$ disebut
 \index{kontraktibel!aljabar-$C^*$}%
\df{kontraktibel} jika ekuivalen secara homotopi dengan~$\{\vc 0\}$.
\end{defn}

\begin{prop}\label{0062327} Misalkan $A$ dan $B$ aljabar-$C^*$ beridentitas. Jika $\phi \sim_h \psi$
dalam $\Hom(A,B)$, maka $K_0(\phi) = K_0(\psi)$.
\end{prop}

%\begin{proof} Lihat \cite{RordamLL:2000}, proposisi 3.2.6.  \ns \end{proof}

\begin{prop}\label{0062331} Jika aljabar-$C^*$ beridentitas $A$ dan $B$
 \index{homotopi!invariansi $K_0$}%
 \index{K0@$K_0$!invariansi homotopi dari}%
ekuivalen secara homotopi, maka $K_0(A) \cong K_0(B)$.
\end{prop}

\begin{prop}\label{0062334} Jika $A$ suatu aljabar-$C^*$ beridentitas, maka barisan eksak terbelah
   \begin{equation}\label{0062334i}
      \vc 0 \to A \to^\iota \wt A \two/->`<-/^Q_\psi \C \to \vc 0
   \end{equation}
\emph{(lihat \ref{0015213}, Kasus 1, butir (7)\,)} menginduksi barisan eksak terbelah lainnya
   \begin{equation}\label{0062334ii}
      \vc 0 \to K_0(A) \to^{K_0(\iota)} K_0(\wt A) \two/->`<-/^{K_0(Q)}_{K_0(\psi)} K_0(\C) \to \vc 0 \,.
   \end{equation}
\end{prop}

\begin{proof}[\emph{Petunjuk untuk bukti}] Definisikan dua fungsi tambahan
   \[ \mu\colon \wt A \sto A\colon a + \lambda \vc j \mapsto a \]
dan
   \[ \psi' \colon \C \sto \wt A \colon \lambda \mapsto \lambda \vc j\,. \]
Kemudian periksa bahwa
   \begin{enumerate}[label=\rm{(\alph*)}]
     \item $\mu \circ \iota = \id A$,
     \item $\iota \circ \mu + \psi' \circ Q = \id{\wt A}$,
     \item $Q \circ \iota = \vc 0$, dan
     \item $Q \circ \psi' = \id{\C}$. % SOURCE-CORRECTION: CH17-C013
   \end{enumerate}   \ns
\end{proof}

\begin{exam} Untuk setiap $n \in \N$,
 \index{K0@$K_0(A)$!adalah $\Z$ ketika $A = \mathbf M_n$}%
$K_0(\mathbf M_n) \cong \Z$.
\end{exam}

\begin{exam} Jika $H$ ruang Hilbert terpisahkan berdimensi tak hingga,
 \index{K0@$K_0(A)$!adalah $\vc 0$ ketika $A = \ofml B(H)$}%
maka $K_0(\ofml B(H)) \cong \vc 0$.
\end{exam}

\begin{defn} Ingat bahwa suatu ruang topologis $X$ disebut
 \index{kontraktibel!ruang topologis}%
\df{kontraktibel} jika terdapat titik $a$ dalam ruang tersebut dan fungsi kontinu $f\colon [0,1] \times X \sto X$ sedemikian sehingga $f(1,x) = x$
dan $f(0,x) = a$ untuk setiap $x \in X$.
\end{defn}

\begin{exam} Jika $X$ ruang Hausdorff kompak yang kontraktibel, maka
 \index{K0@$K_0(A)$!adalah $\Z$ ketika $A = \fml C(X)$, $X$ kontraktibel}%
$K_0(\fml C(X)) \cong \Z$.
\end{exam}


















\section{$\mathbf{\emph{K}_0}(A)$---Kasus Tak Beridentitas}
\begin{defn} Misalkan $A$ aljabar-$C^*$ tak beridentitas. Ingat bahwa barisan eksak terbelah
 ~\eqref{0062334i} bagi unitalisasi $A$ menginduksi barisan eksak terbelah
~\eqref{0062334ii} di antara grup-grup $K_0$ yang bersesuaian.
 \index{K0@$K_0(A)$!untuk $A$ tak beridentitas}%
Definisikan
   \[ \pi:=Q,\qquad K_0(A) = \ker(K_0(\pi))\,. \] % SOURCE-CORRECTION: CH17-C014
\end{defn}

\begin{prop}\label{0062413} Untuk aljabar-$C^*$ tak beridentitas $A$, pemetaan $[\hphantom{P}] \colon \fml
P_\infty(A) \sto K_0(\wt A)$ dapat dipandang sebagai pemetaan dari $\fml P_\infty(A)$ ke dalam~$K_0(A)$.
\end{prop}

\begin{prop}\label{0062416} Untuk aljabar-$C^*$ beridentitas maupun tak beridentitas, barisan
  \[ \vc 0 \to K_0(A) \to K_0(\wt A) \to K_0(\C) \to \vc 0 \]
bersifat eksak.
\end{prop}

\begin{prop}\label{0062418} Untuk aljabar-$C^*$ beridentitas maupun tak beridentitas, grup $K_0(A)$ (isomorfik
dengan) $\ker(K_0(\pi))$.
\end{prop}

\begin{prop}\label{0062424} Jika $\phi\colon A \sto B$ merupakan homomorfisme-$*\,$ di antara aljabar-$C^*$, maka
terdapat homomorfisme grup tunggal % SOURCE-CORRECTION: CH17-C015
 \index{K0@$K_0(\phi)$!kasus tak beridentitas}%
$K_0(\phi)$ yang membuat diagram berikut komutatif.
   \[\xy
     \hSquares/>`>`.>`>`=`>`>/[K_0(A)`K_0(\wt A)`K_0(\C)`K_0(B)`K_0(\wt
     B)`K_0(\C);`K_0(\pi_A)`K_0(\phi)`K_0(\wt\phi)```K_0(\pi_B)]
   \endxy\]
\end{prop}

\begin{prop}\label{0062511} Pasangan pemetaan $A \mapsto K_0(A)$, $\phi \mapsto K_0(\phi)$ merupakan
 \index{funktor!$K_0$}%
 \index{K0@$K_0$!sebagai funktor kovarian}%
funktor kovarian dari kategori $\cat{CSA}$ aljabar-$C^*$ dan homomorfisme-$*\,$ ke
kategori grup Abelian dan homomorfisme grup.
\end{prop}

Dalam proposisi~\ref{0062327} dan~\ref{0062331}, kita menyatakan invariansi homotopi dari
funktor $K_0$ untuk aljabar-$C^*$ beridentitas. Sekarang kita memperluas hasil tersebut ke sebarang
aljabar-$C^*$.

\begin{prop}\label{0062521} Misalkan $A$ dan $B$ aljabar-$C^*$. Jika $\phi \sim_h \psi$
dalam $\Hom(A,B)$, maka $K_0(\phi) = K_0(\psi)$.
\end{prop}

%\begin{proof} Lihat \cite{RordamLL:2000}, proposisi 4.1.4.  \ns \end{proof}

\begin{prop}\label{0062524} Jika aljabar-$C^*$ $A$ dan $B$ ekuivalen secara homotopi,
maka $K_0(A) \cong K_0(B)$.
\end{prop}

%\begin{proof} Lihat \cite{RordamLL:2000}, proposisi 4.1.4.  \ns \end{proof}

\begin{defn} Misalkan $\pi:=Q$ dan $\lambda:=\psi$ adalah homomorfisme-$*\,$ yang berturut-turut berperan sebagai pemetaan hasil bagi dan penampang kanonik dalam barisan eksak
terbelah~\eqref{0062334i} bagi unitalisasi suatu aljabar-$C^*$~$A$. Definisikan % SOURCE-CORRECTION: CH17-C016
 \index{skalar!pemetaan}%
\df{pemetaan skalar} $s\colon \wt A \sto \wt A$ untuk $\wt A$ dengan $s :=  \lambda \circ \pi$. Setiap
anggota $\wt A$ dapat ditulis dalam bentuk $a + \alpha \vc 1_{\wt A}$ untuk suatu $a \in A$ dan
$\alpha \in \C$. Perhatikan bahwa $s(a + \alpha \vc 1_{\wt A}) = \alpha \vc 1_{\wt A}$ dan bahwa $x
- s(x) \in A$ untuk setiap $x \in \wt A$. Untuk setiap bilangan asli $n$, pemetaan skalar menginduksi
pemetaan yang bersesuaian $s = s_n\colon \M n{\wt A} \sto \M n{\wt A}$. Suatu elemen $x \in \M n{\wt
A}$ merupakan
 \index{skalar!elemen}%
\df{elemen skalar} dari $\M n{\wt A}$ jika $s(x) = x$.
\end{defn}

\begin{prop}[Gambaran Standar $K_0(A)$ untuk sebarang $A$]\label{0062544} Jika $A$ suatu
aljabar-$C^*$,
 \index{standar!gambaran $K_0(A)$!untuk sebarang $A$}%
 \index{gambaran (\seeonly{gambaran standar}))}%
maka
     \[ K_0(A) = \{\,[p\,] - [s(p)]\colon p \in \fml P_\infty(\wt A)\}\,. \] % SOURCE-CORRECTION: CH17-C017
\end{prop}























\section{Sifat Keeksakan dan Stabilitas Funktor $K_0$}
\begin{defn} Funktor kovarian $F$ dari kategori $\cat A$ ke kategori $\cat B$ disebut
 \index{eksak terbelah!funktor}%
 \index{funktor!eksak terbelah}%
\df{eksak terbelah} jika membawa barisan eksak terbelah ke barisan eksak terbelah. Funktor tersebut disebut
 \index{funktor eksak separuh}%
 \index{funktor!eksak separuh}%
\df{eksak separuh} jika setiap kali barisan
   \[ \vc 0 \to A_1 \to^j A_2 \to^k A_3 \to \vc 0 \]
eksak dalam $\cat A$, maka
   \[ F(A_1) \to^{F(j)} F(A_2) \to^{F(k)} F(A_3) \]
eksak dalam~$\cat B$.
\end{defn}

\begin{prop}\label{0062664} Funktor $K_0$ eksak separuh.
\end{prop}

%\begin{proof} Lihat~\cite{RordamLL:2000}, proposisi 4.3.2.  \ns \end{proof}

\begin{prop}\label{0062667} Funktor $K_0$ eksak terbelah.
\end{prop}

%\begin{proof} Lihat~\cite{RordamLL:2000}, proposisi 4.3.3.  \ns \end{proof}

\begin{prop}\label{0062671} Funktor $K_0$ mempertahankan jumlah langsung. Artinya, jika $A$ dan $B$
aljabar-$C^*$, maka $K_0(A \oplus B) = K_0(A) \oplus K_0(B)$.
\end{prop}

\begin{exam}\label{0062674} Jika $A$ aljabar-$C^*$, maka $K_0(\wt A) = K_0(A) \oplus \Z$.
\end{exam}

Meskipun eksak terbelah sekaligus eksak separuh, funktor $K_0$ tidak eksak. Masing-masing dari dua
contoh berikut cukup untuk menunjukkan hal ini.

\begin{exam}\label{0062677} Barisan
   \[ \vc 0 \to \fml C_0\bigl((0,1)\bigr) \to^\iota \fml C\bigl([0,1]\bigr)
                               \to^\psi \C \oplus \C \to \vc 0 \]
dengan $\psi(f) = \bigl(f(0),f(1)\bigr)$ jelas eksak; tetapi $K_0(\psi)$ tidak surjektif.
\end{exam}

\begin{exam}\label{0062679} Jika $H$ ruang Hilbert berdimensi tak hingga, barisan eksak % SOURCE-CORRECTION: CH17-C018
  \[ \vc 0 \to \ofml K(H) \to^\iota \ofml B(H) \to^\pi \ofml Q(H) \to \vc 0 \]
yang terkait dengan aljabar Calkin $\ofml Q(H)$ bersifat eksak, tetapi $K_0(\iota)$ tidak injektif.
(Contoh ini memerlukan fakta yang belum kita turunkan: $K_0(\ofml K(H)) \cong \Z$; untuk itu, lihat~\cite{RordamLL:2000}, Akibat~6.4.2.)
\end{exam}

Berikutnya adalah sifat stabilitas penting dari funktor~$K_0$.
\begin{prop}\label{0062681} Jika $A$ aljabar-$C^*$, maka $K_0(A) \cong K_0(\M nA)$.
\end{prop}

%\begin{proof} Lihat~\cite{RordamLL:2000}, proposisi 4.3.8.   \ns \end{proof}





















\section{Limit Induktif}
Dalam bagian~\ref{sec_ind_limits}, kita mendefinisikan \emph{limit induktif} untuk sebarang kategori, dan dalam
bagian~\ref{sec_LFsps}, kita membatasi perhatian pada limit induktif ketat dari ruang vektor topologis
konveks lokal, yang berguna dalam konteks teori distribusi. Dalam bagian
ini, kita menelaah limit induktif dari barisan aljabar-$C^*$.

\begin{defn} Dalam sebarang kategori, suatu
 \index{induktif!barisan}%
 \index{barisan!induktif}%
\df{barisan induktif} adalah pasangan $(A,\phi)$, dengan $A = (A_j)$ barisan objek dan
$\phi = (\phi_j)$ barisan morfisme sedemikian sehingga $\phi_j\colon A_j \sto A_{j+1}$ untuk
setiap~$j$. Suatu
 \index{induktif!limit}%
 \index{limit!induktif}%
\df{limit induktif} (atau
 \index{langsung!limit}%
 \index{limit!langsung}%
\df{limit langsung}) dari barisan $(A,\phi)$ adalah pasangan $(L,\mu)$, dengan $L$ suatu objek dan
$\mu = (\mu_n)$ barisan morfisme $\mu_j\colon A_j \sto L$ yang memenuhi
 \begin{enumerate}
  \item[(i)] $\mu_j = \mu_{j+1} \circ \phi_j$ untuk setiap $j \in \N$, dan
  \item[(ii)] jika $(M,\lambda)$ adalah pasangan dengan $M$ suatu objek dan $\lambda = (\lambda_j)$
suatu barisan morfisme $\lambda_j\colon A_j \sto M$ yang memenuhi $\lambda_j = \lambda_{j+1}
\circ \phi_j$, maka terdapat morfisme tunggal $\psi\colon L \sto M$ sedemikian sehingga $\lambda_j
= \psi \circ \mu_j$ untuk setiap $j \in \N$.
 \end{enumerate}
Dengan sedikit penyalahgunaan bahasa, biasanya kita mengatakan bahwa $L$ adalah limit induktif dari barisan $(A_j)$ dan
 \index{limit@$\underrightarrow{\lim} A_j$ (limit induktif)}%
menulis $L = \underrightarrow{\lim} A_j$.

\vskip 15 pt

   \[\xymatrix@+15pt@C+25pt{&&&& L\ar@{-->}[dddd]^\psi \\
                   &&&&   \\
                  {\cdots}\ar[r] & A_j \ar[uurrr]^{\mu_j}\ar[ddrrr]^{\lambda_j}\ar[r]^(.6){\phi_j} &
                       A_{j+1}\ar[uurr]^{\mu_{j+1}}\ar[ddrr]^{\lambda_{j+1}}\ar[r] & {\cdots} &  \\
                   &&&&   \\
                   &&&& M
    }\]
\end{defn}

\begin{prop} Limit induktif dari barisan induktif (jika ada dalam suatu kategori) tunggal (hingga isomorfisme).
\end{prop}

\begin{prop} Setiap barisan induktif $(A,\phi)$ dari aljabar-$C^*$
 \index{cstaralg@aljabar-$C^*$!limit induktif}%
memiliki limit induktif. (Demikian pula setiap barisan induktif dari grup Abelian.)
\end{prop}

\begin{proof} Lihat \cite{Blackadar:2006}, II.8.2.1; \cite{RordamLL:2000}, proposisi 6.2.4; atau
\cite{Wegge-Olsen:1993}, lampiran~L.  \ns
\end{proof}

\begin{prop} Jika $(L,\mu)$ merupakan limit induktif dari barisan induktif $(A,\phi)$ dari
aljabar-$C^*$, maka $L = \clo{\bigcup {\mu_n}^\sto(A_n)}$ dan $\norm{\mu_m(a)} = \lim_{n \sto
\infty}\norm{\phi_{n,m}(a)} = \inf_{n \ge m}\norm{\phi_{n,m}(a)}$ untuk semua $m \in \N$ dan $a
\in A$.
\end{prop}

\begin{proof} Lihat \cite{RordamLL:2000}, proposisi 6.2.4.   \ns \end{proof}

\begin{defn} Suatu
 \index{aljabar-$C^*$ berdimensi hingga secara aproksimatif}%
 \index{aljabar-AF}%
 \index{hingga!berdimensi!secara aproksimatif}%
 \index{cstaralg@aljabar-$C^*$!aljabar-AF}%
\df{aljabar-$C^*$ berdimensi hingga secara aproksimatif} (singkatnya aljabar-AF) adalah limit induktif
dari suatu barisan aljabar-$C^*$ berdimensi hingga.
\end{defn}

\begin{exam} Jika $(A_n)$ merupakan barisan meningkat yang terdiri atas subaljabar-$C^*$ dari suatu aljabar-$C^*$ $D$
(dengan $\iota_n\colon A_n \sto A_{n+1}$ pemetaan inklusi untuk setiap~$n$), maka $(A,\iota)$
merupakan barisan induktif aljabar-$C^*$ yang limit induktifnya adalah $(B,j)$, dengan $B =
\clo{\bigcup_{n=1}^\infty A_n}$ dan $j_n\colon A_n \sto B$ pemetaan inklusi untuk setiap~$n$. % SOURCE-CORRECTION: CH17-C019
\end{exam}

\begin{exam} Barisan $\mathbf M_1 = \C \to^{\phi_1} \mathbf M_2 \to^{\phi_2} \mathbf M_3
\to^{\phi_3} \dots$ (dengan $\phi_n(a) = \diag(a,0)$ untuk setiap $n \in \N$ dan $a \in \mathbf
M_n$) merupakan barisan induktif aljabar-$C^*$ yang limit induktifnya adalah aljabar-$C^*$
$\ofml K(H)$ yang terdiri atas operator kompak pada ruang Hilbert terpisahkan berdimensi tak hingga~$H$. % SOURCE-CORRECTION: CH17-C020
\end{exam}

\begin{proof} Lihat \cite{RordamLL:2000}, bagian 6.4.  \ns \end{proof}

\begin{exam} Barisan $\Z \to^{\vc 1} \Z \to^{\vc 2} \Z \to^{\vc 3} \Z \to^{\vc 4} \dots$
(dengan $\vc n \colon \Z \sto \Z$ memenuhi $\vc n(1) = n$) merupakan barisan induktif grup
Abelian yang limit induktifnya adalah himpunan $\Q$ dari bilangan rasional.
\end{exam}

\begin{exam} Barisan $\Z \to^{\vc 2} \Z \to^{\vc 2} \Z \to^{\vc 2} \Z \to^{\vc 2} \dots$
merupakan barisan induktif grup Abelian yang limit induktifnya adalah himpunan bilangan rasional
diadik.
\end{exam}

Hasil berikut disebut
 \index{kekontinuan!dari $K_0$}%
 \index{K0@$K_0$!kekontinuan dari}%
\emph{sifat kekontinuan $K_0$}.

\begin{prop} Jika $(A,\phi)$ barisan induktif aljabar-$C^*$, maka
   \[ K_0\bigl(\,\underrightarrow{\lim}\, A_n\bigr) = \underrightarrow{\lim}\,K_0(A_n)\,. \]
\end{prop}

\begin{proof} Lihat \cite{RordamLL:2000}, teorema 6.3.2.   \ns \end{proof}

\begin{exam} Jika $H$ ruang Hilbert,
 \index{K0@$K_0(A)$!adalah $\Z$ ketika $A = \ofml K(H)$}%
maka $K_0(\ofml K(H)) \cong \Z$.
\end{exam}





















\section{Diagram Bratteli}
\begin{prop} Homomorfisme-$*\,$ tak nol dari $M_k$ ke $M_n$ ada hanya jika $n \ge k$; dalam
hal ini, semuanya tepat merupakan pemetaan berbentuk
   \[ \phi\colon a \mapsto u\,\diag(a, a, \dots, a, \vc 0)\, u^* \] % SOURCE-CORRECTION: CH17-C021
dengan $u$ matriks uniter. Di sini terdapat $m$ salinan $a$, dan $\vc 0$ adalah matriks nol berukuran $r \times
r$, dengan $n = mk + r$. Bilangan $m$ adalah
 \index{multiplisitas!dari homomorfisme-$*\,$}%
\df{multiplisitas} dari~$\phi$.
\end{prop}

\begin{proof} Lihat \cite{Power:1992}, akibat 1.3.  \ns \end{proof}

\begin{exam} Contoh homomorfisme-$*\,$ dari $M_2$ ke $M_7$ adalah
 \[  A =
  \begin{bmatrix} a & b \\ c & d  \end{bmatrix} \mapsto
    \begin{bmatrix}
      a & b & 0 & 0 & 0 & 0 & 0 \\
      c & d & 0 & 0 & 0 & 0 & 0 \\
      0 & 0 & a & b & 0 & 0 & 0 \\
      0 & 0 & c & d & 0 & 0 & 0 \\
      0 & 0 & 0 & 0 & 0 & 0 & 0 \\
      0 & 0 & 0 & 0 & 0 & 0 & 0 \\
      0 & 0 & 0 & 0 & 0 & 0 & 0
    \end{bmatrix} = \diag(A,A, \vc 0)\]
 dengan $m = 2$ dan $r = 3$.
\end{exam}

\enlargethispage{\baselineskip}

\begin{prop}\label{0068117} Setiap aljabar-$C^*$ berdimensi hingga $A$ isomorfik dengan
jumlah langsung aljabar matriks
   \[ A \simeq M_{k_1} \oplus \dots \oplus M_{k_r}.\]
Misalkan $A$ dan $B$ aljabar-$C^*$ berdimensi hingga, sehingga
   \[ A \simeq M_{k_1} \oplus \dots \oplus M_{k_r} \qquad \text{dan}
                          \qquad B \simeq M_{n_1} \oplus \dots \oplus M_{n_s} \]
dan misalkan $\phi\colon A \sto B$ homomorfisme-$*\,$ beridentitas. Maka $\phi$
ditentukan (hingga ekuivalensi uniter dalam~$B$) oleh matriks $s \times r$ $\vc m = \bigl[
m_{ij}\bigr]$ yang berentri bilangan bulat tak negatif sedemikian sehingga % SOURCE-CORRECTION: CH17-C022
   \begin{equation}\label{0068117i}
    \vc m\, \vc k = \vc n.
   \end{equation}
\emph{Di sini $\vc k = (k_1, \dots,k_r)$, $\vc n = (n_1, \dots, n_s)$, dan bilangan $m_{ij}$ adalah
multiplisitas pemetaan}
 \[\xymatrix@1{M_{k_j} \ar[r] & A \ar[r]^\phi & B \ar[r] & M_{n_i}
 }\]
\end{prop}

\begin{proof} Lihat \cite{Davidson:1996}, teorema III.1.1.  \ns \end{proof}

\begin{defn} Misalkan $\phi$ seperti dalam proposisi~\ref{0068117} sebelumnya. Suatu
 \index{diagram Bratteli}%
 \index{diagram!Bratteli}%
\df{diagram Bratteli} untuk $\phi$ terdiri atas dua baris (atau kolom) simpul yang diberi label
$k_j$ dan $n_i$, bersama $m_{ij}$ sisi yang menghubungkan $k_j$ ke $n_i$ (untuk $1 \le j \le
r$ dan $1 \le i \le s$).
\end{defn}

\begin{exam} Misalkan $\phi\colon \C \oplus \C \sto M_4 \oplus M_3$ diberikan oleh
 \[ (\lambda, \mu) \mapsto \left(
      \begin{bmatrix}
        \lambda & 0 & 0 & 0 \\
        0 & \lambda & 0 & 0 \\
        0 & 0 & \lambda & 0 \\
        0 & 0 & 0 & \mu
      \end{bmatrix}\, , \,
      \begin{bmatrix}
        \lambda & 0 & 0 \\
        0 & \lambda & 0 \\
        0 & 0 & \mu
      \end{bmatrix}\right)\, .\]
Maka $m_{11} = 3$, $m_{12} = 1$, $m_{21} = 2$, dan $m_{22} = 1$. Perhatikan bahwa
   \[\vc m \,\vc k = \begin{bmatrix} 3 & 1 \\ 2 & 1 \end{bmatrix}\,
             \begin{bmatrix} 1 \\ 1 \end{bmatrix} =
             \begin{bmatrix} 4 \\ 3 \end{bmatrix} = \vc n.\]
Suatu diagram Bratteli untuk $\phi$ adalah
    \[\xymatrix@=90pt{1 \ar@<.7ex>[r]\ar[r]\ar@<-.7ex>[r]\ar@<.5ex>[dr]\ar@<-.5ex>[dr] & 4 \\
                1 \ar[ur]\ar[r] & 3
     }\]
\end{exam}

\begin{exam} Misalkan $\phi\colon \C \oplus M_2 \sto M_3 \oplus M_5 \oplus M_2$
diberikan oleh
 \[ (\lambda, b) \mapsto \left(
    \begin{bmatrix}
        \lambda & 0  \\
            0 &  b
      \end{bmatrix}\, , \,
      \begin{bmatrix}
        \lambda & 0 & 0  \\
              0 & b & 0  \\
              0 & 0 &  b
      \end{bmatrix} \,,\, b \, \right).
    \]
Maka $m_{11} = 1$, $m_{12} = 1$, $m_{21} = 1$, $m_{22} = 2$, $m_{31} = 0$, dan $m_{32} = 1$.
Perhatikan bahwa
 \[\vc m \,\vc k = \begin{bmatrix} 1 & 1 \\ 1 & 2 \\ 0 & 1 \end{bmatrix}\,
             \begin{bmatrix} 1 \\ 2 \end{bmatrix} =
             \begin{bmatrix} 3 \\ 5 \\2 \end{bmatrix} = \vc n.\]
Suatu diagram Bratteli untuk $\phi$ adalah
 \[\xymatrix@=20pt@C+30pt{&& 3                    \\
                   1 \ar[urr]\ar[drr] &&   \\
                   && 5                    \\
         2\ar[uuurr]\ar@<+.5ex>[urr]\ar@<-.5ex>[urr]\ar[drr] && \\
                   && 2
   }\]
\end{exam}

\begin{exam} Misalkan $\phi\colon M_3 \oplus M_2 \oplus M_2 \sto M_9 \oplus M_7$
diberikan oleh

\vskip 10 pt

 \[ (a,b,c) \mapsto \left(
    \begin{bmatrix}
            a & \vc 0 & \vc 0 \\
        \vc 0 &     a & \vc 0 \\
        \vc 0 & \vc 0 & \vc 0
      \end{bmatrix}\, , \,
      \begin{bmatrix}
            a & \vc 0 & \vc 0  \\
        \vc 0 &     b & \vc 0  \\
        \vc 0 & \vc 0 &     c
      \end{bmatrix}\, \right)\, .\]
Maka $m_{11} = 2$, $m_{12} = 0$, $m_{13} = 0$, $m_{21} = 1$, $m_{22} = 1$, dan $m_{23} = 1$.
Perhatikan bahwa kali ini terjadi sesuatu yang ganjil: kita mempunyai $\vc m\,\vc k \ne \vc n$:
 \[\vc m \,\vc k = \begin{bmatrix} 2 & 0 & 0 \\ 1 & 1 & 1 \end{bmatrix}\,
             \begin{bmatrix} 3 \\ 2 \\ 2 \end{bmatrix} =
             \begin{bmatrix} 6 \\ 7 \end{bmatrix} \le
             \begin{bmatrix} 9 \\ 7 \end{bmatrix} = \vc n.\]
Apa masalahnya di sini? Nah, hasil yang dinyatakan sebelumnya berlaku untuk homomorfisme-$*\,$
\emph{beridentitas}, sedangkan $\phi$ ini \emph{tidak} beridentitas. Secara umum, inilah hasil terbaik
yang dapat kita harapkan: $\vc m\,\vc k \le \vc n$.

\noindent Berikut diagram Bratteli yang dihasilkan untuk $\phi$:
   \[\xymatrix@=+20pt@C+20pt@R-7pt{
        3\ar@<+.5ex>[drr]\ar@<-.5ex>[drr]\ar[dddrr] && \\
                            && 9                   \\
                           2\ar[drr] &&            \\
                            && 7                   \\
                           2\ar[urr] &&
   }\]
\end{exam}

\begin{exam} Jika $K$ himpunan Cantor, maka $\fml C(K)$ merupakan aljabar-AF.
 \index{AF@aljabar-AF!$\fml C(K)$ merupakan aljabar-AF ($K$ himpunan Cantor)}%
Untuk melihatnya, tuliskan $K$ sebagai irisan dari suatu keluarga menurun yang terdiri atas himpunan bagian tertutup $K_j$, yang masing-masing
terdiri atas $2^j$ subinterval tertutup yang saling lepas dari $[0,1]$. Untuk setiap $j \ge 0$, misalkan
$A_j$ subaljabar yang terdiri atas fungsi-fungsi dalam $\fml C(K)$ yang konstan pada setiap
interval penyusun~$K_j$. Jadi $A_j \simeq \C^{2^j}$ untuk setiap~$j \ge 0$. Pembenaman
 \[\phi_j \colon A_j \sto A_{j+1}\colon
      (a_1, a_2, \dots, a_{2^j}) \mapsto
        (a_1, a_1, a_2, a_2,\dots, a_{2^j}, a_{2^j})\]
memecah setiap proyeksi minimal menjadi jumlah dua proyeksi minimal. Untuk $\phi_0$, matriks
$\vc m_0$ dari ``multiplisitas parsial'' yang bersesuaian adalah $\begin{bmatrix} 1 \\ 1
\end{bmatrix}$; matriks $\vc m_1$ yang bersesuaian dengan $\phi_1$ adalah
$\begin{bmatrix} 1 & 0\\ 1 & 0 \\ 0 & 1 \\ 0 & 1 \end{bmatrix}$; dan seterusnya. Jadi diagram Bratteli
untuk limit induktif $\fml C(K) = \underrightarrow{\lim}A_j$ adalah:

\vskip 20 pt

 \[\xymatrix@=+7pt@C+20pt@R-7pt{
     &&& 1 & \\
     && 1\ar[ur]\ar[dr] && \\
     &&& 1 & \\
     & 1\ar[uur]\ar[ddr] &&& \\
     &&& 1 & \\
     && 1\ar[ur]\ar[dr] && \\
     &&& 1 & \\
     1\ar[uuuur]\ar[ddddr] &&&& \cdots \\
     &&& 1 & \\
     && 1\ar[ur]\ar[dr] && \\
     &&& 1 & \\
     & 1\ar[uur]\ar[ddr] &&& \\
     &&&1 & \\
     && 1\ar[ur]\ar[dr] && \\
     &&& 1 &
   }\]
\end{exam}

\begin{exam} Berikut suatu contoh dari apa yang disebut
 \index{aljabar CAR}%
 \index{cstaralg@aljabar-$C^*$!aljabar CAR}%
aljabar CAR (CAR = Relasi Antikomutasi Kanonik). Untuk $j \ge 0$, misalkan $A_j = % SOURCE-CORRECTION: CH17-C023
M_{\spsp{2}j}$ dan
 \[\phi_j\colon A_j \to A_{j+1}\colon \vc a \mapsto
      \begin{bmatrix} \vc a & 0 \\ 0 & \vc a \end{bmatrix}.\]
``Matriks'' multiplisitas $\vc m$ untuk setiap $\phi_j$ hanyalah matriks $1 \times 1$~$[2]$.
Kita melihat bahwa untuk setiap $j$
 \[\vc m\, \vc k(j) = [2]\,[2^j] = [2^{j+1}] = \vc n(j).\]
Ini menghasilkan diagram Bratteli berikut

\vskip 15pt

 \[\xymatrix{1 \ar@<.5ex>[r]\ar@<-.5ex>[r] & 2 \ar@<.5ex>[r]\ar@<-.5ex>[r] & 4
        \ar@<.5ex>[r]\ar@<-.5ex>[r] & 8 \ar@<.5ex>[r]\ar@<-.5ex>[r] & 16
        \ar@<.5ex>[r]\ar@<-.5ex>[r] & 32
        \ar@<.5ex>[r]\ar@<-.5ex>[r]& {\cdots}
   }\]
\vskip 15pt
\noindent untuk limit induktif $C = \underrightarrow{\lim}A_j$.

\textbf{Namun}, jangkauan $\phi_j$ termuat dalam subaljabar $B_{j+1} \simeq M_{2^j}
\oplus M_{2^j}$ dari~$A_{j+1}$. (Ambil $B_0 = A_0 = \C$.) Jadi, untuk $j \in \N$, kita dapat memandang
$\phi_j$ sebagai pemetaan dari $B_j$ ke~$B_{j+1}$:
 \[ \phi_j\colon(\vc b, \vc c) \mapsto
    \left(\begin{bmatrix} \vc b & 0 \\ 0 & \vc c \end{bmatrix},
          \begin{bmatrix} \vc b & 0 \\ 0 & \vc c \end{bmatrix}
              \right)\,.\]
Sekarang matriks multiplisitas $\vc m$ untuk setiap $\phi_j$ adalah $\begin{bmatrix} 1 & 1 \\ 1 & 1
\end{bmatrix}$, dan kita melihat bahwa untuk setiap $j$
 \[\vc m\, \vc k(j) =
   \begin{bmatrix} 1 & 1 \\ 1 & 1 \end{bmatrix}\,
        \begin{bmatrix} 2^{j-1} \\ 2^{j-1} \end{bmatrix} =
   \begin{bmatrix} 2^j \\ 2^j \end{bmatrix} = \vc n(j).\]
Karena $C = \underrightarrow{\lim}A_j = \underrightarrow{\lim}B_j$, kita memperoleh diagram Bratteli kedua
(yang cukup berbeda) untuk aljabar-AF~$C$.
   \[\xymatrix@-10pt{&& 1 \ar[rr]\ar[rrdd] && 2 \ar[rr]\ar[rrdd] && 4 \ar[rr]\ar[rrdd]
                                && 8 \ar[rr]\ar[rrdd] && 16 \ar[rr]\ar[rrdd]&& \cdots\\
              1 \ar[urr]\ar[drr]&&&&&&&&&&  \\
             && 1 \ar[rr]\ar[rruu] && 2 \ar[rr]\ar[rruu] && 4 \ar[rr]\ar[rruu]
                                && 8 \ar[rr]\ar[rruu] && 16 \ar[rr]\ar[rruu]&& \cdots \\
   }\]
\end{exam}

\begin{exam} Ini adalah
 \index{aljabar Fibonacci}%
 \index{cstaralg@aljabar-$C^*$!aljabar Fibonacci}%
\df{aljabar Fibonacci}. Untuk $j \in \N$, definisikan barisan $\bigl(p_j\bigr)$ dan
$\bigl(q_j\bigr)$ dengan relasi rekursi yang dikenal:
 \begin{align*}
       p_1 = &q_1 = 1, \\
       p_{j+1} = &p_j + q_j, \text{ dan } \\
       q_{j+1} &= p_j.
 \end{align*}
Untuk semua $j \in \N$, misalkan $A_j = M_{p_j} \oplus M_{q_j}$ dan

\vskip 10 pt

 \[\phi_j\colon A_j \sto A_{j+1}\colon (a,b) \mapsto
     \left(\begin{bmatrix} a & 0 \\ 0 & b
          \end{bmatrix},  a \right). \]
Matriks multiplisitas $\vc m$ untuk setiap $\phi_j$ adalah $\begin{bmatrix} 1 & 1 \\ 1 & 0
\end{bmatrix}$, dan untuk setiap $j$
 \[\vc m\, \vc k(j) =
   \begin{bmatrix} 1 & 1 \\ 1 & 0 \end{bmatrix}\,
        \begin{bmatrix} p_j \\ q_j \end{bmatrix} =
   \begin{bmatrix} p_j + q_j \\ p_j \end{bmatrix} =
        \begin{bmatrix} p_{j+1} \\ q_{j+1} \end{bmatrix} =
                          \vc n(j).\]
Diagram Bratteli yang dihasilkan untuk $F = \underrightarrow{\lim}A_j$ adalah
   \[\xymatrix@-10pt{&& 1 \ar[rr]\ar[rrdd] && 2 \ar[rr]\ar[rrdd] &&
         3 \ar[rr]\ar[rrdd] && 5 \ar[rr]\ar[rrdd] && 8 \ar[rr]\ar[rrdd]&& \cdots\\
              1 \ar[urr]\ar[drr]&&&&&&&&&&  \\
             && 1 \ar[rruu] && 1 \ar[rruu] && 2 \ar[rruu] &&
              3 \ar[rruu] && 5 \ar[rruu]&& \cdots \\
   }\]
\end{exam}













\endinput
